\documentclass[10pt,a4paper]{article}
\usepackage{amsmath,amssymb,amsthm,graphicx,verbatim,fontspec,hyperref}
\setmainfont{lmroman10-regular.otf}
% Retained preamble for the Stacks Project chapter snapshot in this directory.
% The chapter's own preamble.tex is a copy of the shared Stacks Project preamble; the block
% below is sliced MECHANICALLY, line for line, from it. Nothing was added, removed or
% redefined; the wrapper class, the xr-hyper/hyperref wiring and the externaldocument list
% are omitted because the assembler supplies the document class.
%
% lines taken from preamble.tex: 3, 5, 9, 11, 12, 156, 157, 158, ... (51 lines total)
%
% package coverage restored from the chapter's own preamble (2026-10-13): xy, xy-2cell

\usepackage[all]{xy}
\xyoption{2cell}
\UseAllTwocells

\newenvironment{reference}{\comment}{\endcomment}
\newenvironment{slogan}{\comment}{\endcomment}
\newenvironment{history}{\comment}{\endcomment}
\theoremstyle{plain}
\newtheorem{theorem}[subsection]{Theorem}
\newtheorem{proposition}[subsection]{Proposition}
\newtheorem{lemma}[subsection]{Lemma}
\theoremstyle{definition}
\newtheorem{definition}[subsection]{Definition}
\newtheorem{example}[subsection]{Example}
\newtheorem{exercise}[subsection]{Exercise}
\newtheorem{situation}[subsection]{Situation}
\theoremstyle{remark}
\newtheorem{remark}[subsection]{Remark}
\newtheorem{remarks}[subsection]{Remarks}
\numberwithin{equation}{subsection}
\def\lim{\mathop{\mathrm{lim}}\nolimits}
\def\colim{\mathop{\mathrm{colim}}\nolimits}
\def\Spec{\mathop{\mathrm{Spec}}}
\def\Hom{\mathop{\mathrm{Hom}}\nolimits}
\def\Ext{\mathop{\mathrm{Ext}}\nolimits}
\def\SheafHom{\mathop{\mathcal{H}\!\mathit{om}}\nolimits}
\def\SheafExt{\mathop{\mathcal{E}\!\mathit{xt}}\nolimits}
\def\Sch{\mathit{Sch}}
\def\Mor{\mathop{\mathrm{Mor}}\nolimits}
\def\Ob{\mathop{\mathrm{Ob}}\nolimits}
\def\Sh{\mathop{\mathit{Sh}}\nolimits}
\def\NL{\mathop{N\!L}\nolimits}
\def\CH{\mathop{\mathrm{CH}}\nolimits}
\def\proetale{{pro\text{-}\acute{e}tale}}
\def\etale{{\acute{e}tale}}
\def\QCoh{\mathit{QCoh}}
\def\Ker{\mathop{\mathrm{Ker}}}
\def\Im{\mathop{\mathrm{Im}}}
\def\Coker{\mathop{\mathrm{Coker}}}
\def\Coim{\mathop{\mathrm{Coim}}}
\DeclareMathSymbol{\boxtimes}{\mathbin}{AMSa}{"02}
\def\QCohstack{\mathcal{QC}\!\mathit{oh}}
\def\Cohstack{\mathcal{C}\!\mathit{oh}}
\def\Spacesstack{\mathcal{S}\!\mathit{paces}}
\def\Quotfunctor{\mathrm{Quot}}
\def\Hilbfunctor{\mathrm{Hilb}}
\def\Curvesstack{\mathcal{C}\!\mathit{urves}}
\def\Polarizedstack{\mathcal{P}\!\mathit{olarized}}
\def\Complexesstack{\mathcal{C}\!\mathit{omplexes}}
\def\Pic{\mathop{\mathrm{Pic}}\nolimits}
\def\Picardstack{\mathcal{P}\!\mathit{ic}}
\def\Picardfunctor{\mathrm{Pic}}
\def\Deformationcategory{\mathcal{D}\!\mathit{ef}}
\title{Brown representability under detecting and countably factorizing generators}
\author{The Stacks Project Authors}
\date{stacks-project:derived.tex@a04446e57ec1fbc252a871afcec7752fb2807b14; GNU FDL 1.2 or later}
\begin{document}
\maketitle
\noindent\textit{Editorial note.} Original theorem, complete author proof, surrounding source sections and shared bibliography retained. Cross-section/chapter references outside the collected sections display their exact original labels, without invented theorem numbers. No hypothesis or mathematical argument added or repaired.\par\medskip
\noindent\textit{Allowed standard prerequisites (editorial).}\begin{itemize}
\item Additive categories, categories of fractions, triangulated categories and basic homological algebra. All selected theorem constructions and the surrounding local definitions and proofs are included. Standard external references retain their exact original labels and shared bibliography.. No primary localization or representing-object construction is assumed as background.
\end{itemize}\medskip
\section{Derived colimits}
\label{section-derived-colimit}

\noindent
In a triangulated category there is a notion of derived colimit.

\begin{definition}
\label{definition-derived-colimit}
Let $\mathcal{D}$ be a triangulated category.
Let $(K_n, f_n)$ be a system of objects of $\mathcal{D}$.
We say an object $K$ is a {\it derived colimit}, or a
{\it homotopy colimit} of the system $(K_n)$ if
the direct sum $\bigoplus K_n$ exists and there is a distinguished triangle
$$
\bigoplus K_n \to \bigoplus K_n \to K \to \bigoplus K_n[1]
$$
where the map $\bigoplus K_n \to \bigoplus K_n$ is given
by $1 - f_n$ in degree $n$. If this is the
case, then we sometimes indicate this by the notation
$K = \text{hocolim} K_n$.
\end{definition}

\noindent
By TR3 a derived colimit, if it exists, is unique up to (non-unique)
isomorphism. Moreover, by TR1 a derived colimit of $K_n$ exists
as soon as $\bigoplus K_n$ exists. The derived category $D(\textit{Ab})$
of the category of abelian groups is an example of a triangulated
category where all homotopy colimits exist.

\medskip\noindent
The nonuniqueness makes it hard to pin down the derived colimit. In
More on Algebra, Lemma \textnormal{[original label: \nolinkurl{more-algebra-lemma-map-from-hocolim}]}
the reader finds an exact sequence
$$
0 \to R^1\lim \Hom(K_n, L[-1]) \to \Hom(\text{hocolim} K_n, L)
\to \lim \Hom(K_n, L) \to 0
$$
describing the $\Hom$s out of a homotopy colimit in terms of the
usual $\Hom$s.

\begin{remark}
\label{remark-uniqueness-derived-colimit}
Let $\mathcal{D}$ be a triangulated category.
Let $(K_n, f_n)$ be a system of objects of $\mathcal{D}$.
We may think of a derived colimit as an object $K$
of $\mathcal{D}$ endowed with morphisms $i_n : K_n \to K$
such that $i_{n + 1} \circ f_n = i_n$ and such that there
exists a morphism $c : K \to \bigoplus K_n[1]$ with the property that
$$
\bigoplus K_n \xrightarrow{1 - f_n} \bigoplus K_n \xrightarrow{i_n}
K \xrightarrow{c} \bigoplus K_n[1]
$$
is a distinguished triangle. If $(K', i'_n, c')$ is a second
derived colimit, then there exists an isomorphism
$\varphi : K \to K'$ such that $\varphi \circ i_n = i'_n$ and
$c' \circ \varphi = c$. The existence of $\varphi$ is
TR3 and the fact that $\varphi$ is an isomorphism is
Lemma \textnormal{[original label: \nolinkurl{lemma-third-isomorphism-triangle}]}.
\end{remark}

\begin{remark}
\label{remark-functoriality-derived-colimit}
Let $\mathcal{D}$ be a triangulated category.
Let $(a_n) : (K_n, f_n) \to (L_n, g_n)$ be a morphism of systems
of objects of $\mathcal{D}$. Let $(K, i_n, c)$ be a derived
colimit of the first system and let $(L, j_n, d)$ be a derived
colimit of the second system with notation as in
Remark \ref{remark-uniqueness-derived-colimit}.
Then there exists a morphism $a : K \to L$
such that $a \circ i_n = j_n$ and $d \circ a = (a_n[1]) \circ c$.
This follows from TR3 applied to the defining distinguished
triangles.
\end{remark}

\begin{lemma}
\label{lemma-hocolim-subsequence}
Let $\mathcal{D}$ be a triangulated category.
Let $(K_n, f_n)$ be a system of objects of $\mathcal{D}$.
Let $n_1 < n_2 < n_3 < \ldots$ be a sequence of integers.
Assume $\bigoplus K_n$ and $\bigoplus K_{n_i}$ exist.
Then there exists an isomorphism
$\text{hocolim} K_{n_i} \to \text{hocolim} K_n$
such that
$$
\xymatrix{
K_{n_i} \ar[r] \ar[d]_{\text{id}} & \text{hocolim} K_{n_i} \ar[d] \\
K_{n_i} \ar[r] & \text{hocolim} K_n
}
$$
commutes for all $i$.
\end{lemma}

\begin{proof}
Let $g_i : K_{n_i} \to K_{n_{i + 1}}$ be the composition
$f_{n_{i + 1} - 1} \circ \ldots \circ f_{n_i}$.
We construct commutative diagrams
$$
\vcenter{
\xymatrix{
\bigoplus\nolimits_i K_{n_i} \ar[r]_{1 - g_i} \ar[d]_b &
\bigoplus\nolimits_i K_{n_i} \ar[d]^a \\
\bigoplus\nolimits_n K_n \ar[r]^{1 - f_n} &
\bigoplus\nolimits_n K_n
}
}
\quad\text{and}\quad
\vcenter{
\xymatrix{
\bigoplus\nolimits_n K_n \ar[r]_{1 - f_n} \ar[d]_d &
\bigoplus\nolimits_n K_n \ar[d]^c \\
\bigoplus\nolimits_i K_{n_i} \ar[r]^{1 - g_i} &
\bigoplus\nolimits_i K_{n_i}
}
}
$$
as follows. Let $a_i = a|_{K_{n_i}}$ be the inclusion of $K_{n_i}$
into the direct sum. In other words, $a$ is the natural inclusion.
Let $b_i = b|_{K_{n_i}}$ be the map
$$
K_{n_i}
\xrightarrow{1,\ f_{n_i},\ f_{n_i + 1} \circ f_{n_i},
\ \ldots,\ f_{n_{i + 1} - 2} \circ \ldots \circ f_{n_i}}
K_{n_i} \oplus K_{n_i + 1} \oplus \ldots \oplus K_{n_{i + 1} - 1}
$$
If $n_{i - 1} < j \leq n_i$, then we let $c_j = c|_{K_j}$
be the map
$$
K_j \xrightarrow{f_{n_i - 1} \circ \ldots \circ f_j} K_{n_i}
$$
We let $d_j = d|_{K_j}$ be zero if $j \not = n_i$ for any $i$
and we let $d_{n_i}$ be the natural inclusion of $K_{n_i}$
into the direct sum. In other words, $d$ is the natural projection.
By TR3 these diagrams define morphisms
$$
\varphi : \text{hocolim} K_{n_i} \to \text{hocolim} K_n
\quad\text{and}\quad
\psi : \text{hocolim} K_n \to \text{hocolim} K_{n_i}
$$
Since $c \circ a$ and $d \circ b$ are the identity maps we see that
$\varphi \circ \psi$ is an isomorphism by
Lemma \textnormal{[original label: \nolinkurl{lemma-third-isomorphism-triangle}]}.
The other way around we get the morphisms $a \circ c$ and $b \circ d$.
Consider the morphism
$h = (h_j) : \bigoplus K_n \to \bigoplus K_n$ given by
the rule: for $n_{i - 1} < j < n_i$ we set
$$
h_j : K_j
\xrightarrow{1,\ f_j,\ f_{j + 1} \circ f_j,
\ \ldots,\ f_{n_i - 1} \circ \ldots \circ f_j}
K_j \oplus \ldots \oplus K_{n_i}
$$
Then the reader verifies that $(1 - f) \circ h = \text{id} - a \circ c$
and $h \circ (1 - f) = \text{id} - b \circ d$. This means that
$\text{id} - \psi \circ \varphi$ has square zero by
Lemma \textnormal{[original label: \nolinkurl{lemma-third-map-square-zero}]} (small argument omitted).
In other words, $\psi \circ \varphi$ differs from the identity
by a nilpotent endomorphism, hence is an isomorphism. Thus
$\varphi$ and $\psi$ are isomorphisms as desired.
\end{proof}

\begin{lemma}
\label{lemma-direct-sums}
Let $\mathcal{A}$ be an abelian category.
If $\mathcal{A}$ has exact countable direct sums, then
$D(\mathcal{A})$ has countable direct sums. In fact given
a collection of complexes $K_i^\bullet$ indexed by a countable
index set $I$ the termwise direct sum $\bigoplus K_i^\bullet$
is the direct sum of $K_i^\bullet$ in $D(\mathcal{A})$.
\end{lemma}

\begin{proof}
Let $L^\bullet$ be a complex. Suppose given maps
$\alpha_i : K_i^\bullet \to L^\bullet$ in $D(\mathcal{A})$.
This means there exist quasi-isomorphisms
$s_i : M_i^\bullet \to K_i^\bullet$
of complexes and maps of complexes $f_i : M_i^\bullet \to L^\bullet$
such that $\alpha_i = f_is_i^{-1}$. By assumption the map of complexes
$$
s : \bigoplus M_i^\bullet \longrightarrow \bigoplus K_i^\bullet
$$
is a quasi-isomorphism. Hence setting $f = \bigoplus f_i$ we see that
$\alpha = fs^{-1}$ is a map in $D(\mathcal{A})$ whose composition
with the coprojection $K_i^\bullet \to \bigoplus K_i^\bullet$ is $\alpha_i$.
We omit the verification that $\alpha$ is unique.
\end{proof}

\begin{lemma}
\label{lemma-compute-colimit}
Let $\mathcal{A}$ be an abelian category. Assume colimits over $\mathbf{N}$
exist and are exact. Then countable direct sums exists and are exact.
Moreover, if $(A_n, f_n)$ is a system over $\mathbf{N}$, then there is
a short exact sequence
$$
0 \to \bigoplus A_n \to \bigoplus A_n \to \colim A_n \to 0
$$
where the first map in degree $n$ is given by $1 - f_n$.
\end{lemma}

\begin{proof}
The first statement follows from
$\bigoplus A_n = \colim (A_1 \oplus \ldots \oplus A_n)$.
For the second, note that for each $n$ we have the short exact sequence
$$
0 \to
A_1 \oplus \ldots \oplus A_{n - 1} \to
A_1 \oplus \ldots \oplus A_n \to A_n \to 0
$$
where the first map is given by the maps $1 - f_i$ and the
second map is the sum of the transition maps.
Take the colimit to get the sequence of the lemma.
\end{proof}

\begin{lemma}
\label{lemma-colim-hocolim}
Let $\mathcal{A}$ be an abelian category. Let $L_n^\bullet$
be a system of complexes of $\mathcal{A}$. Assume
colimits over $\mathbf{N}$ exist and are exact in $\mathcal{A}$.
Then the termwise
colimit $L^\bullet = \colim L_n^\bullet$ is a homotopy colimit of the
system in $D(\mathcal{A})$.
\end{lemma}

\begin{proof}
We have an exact sequence of complexes
$$
0 \to \bigoplus L_n^\bullet \to \bigoplus L_n^\bullet \to L^\bullet \to 0
$$
by Lemma \ref{lemma-compute-colimit}.
The direct sums are direct sums in $D(\mathcal{A})$ by
Lemma \ref{lemma-direct-sums}.
Thus the result follows from the definition
of derived colimits in
Definition \ref{definition-derived-colimit}
and the fact that a short exact sequence of complexes
gives a distinguished triangle
(Lemma \textnormal{[original label: \nolinkurl{lemma-derived-canonical-delta-functor}]}).
\end{proof}

\begin{lemma}
\label{lemma-cohomology-of-hocolim}
Let $\mathcal{D}$ be a triangulated category having countable
direct sums. Let $\mathcal{A}$ be an abelian category with exact
colimits over $\mathbf{N}$.
Let $H : \mathcal{D} \to \mathcal{A}$ be a homological functor
commuting with countable direct sums.
Then $H(\text{hocolim} K_n) = \colim H(K_n)$
for any system of objects of $\mathcal{D}$.
\end{lemma}

\begin{proof}
Write $K = \text{hocolim} K_n$. Apply $H$ to the defining
distinguished triangle to get
$$
\bigoplus H(K_n) \to \bigoplus H(K_n)
\to H(K) \to
\bigoplus H(K_n[1]) \to \bigoplus H(K_n[1])
$$
where the first map is given by $1 - H(f_n)$ and the last map
is given by $1 - H(f_n[1])$.
Apply Lemma \ref{lemma-compute-colimit} to see that this proves the lemma.
\end{proof}

\noindent
The following lemma tells us that taking maps out of a compact
object (to be defined later) commutes with derived colimits.

\begin{lemma}
\label{lemma-commutes-with-countable-sums}
Let $\mathcal{D}$ be a triangulated category with countable direct sums.
Let $K \in \mathcal{D}$ be an object such that for every
countable set of objects $E_n \in \mathcal{D}$ the canonical map
$$
\bigoplus \Hom_\mathcal{D}(K, E_n)
\longrightarrow
\Hom_\mathcal{D}(K, \bigoplus E_n)
$$
is a bijection. Then, given any system $L_n$ of $\mathcal{D}$ over
$\mathbf{N}$ whose derived colimit $L = \text{hocolim} L_n$
exists we have that
$$
\colim \Hom_\mathcal{D}(K, L_n) \longrightarrow \Hom_\mathcal{D}(K, L)
$$
is a bijection.
\end{lemma}

\begin{proof}
By Lemma \textnormal{[original label: \nolinkurl{lemma-representable-homological}]}, this is a particular case
of Lemma \ref{lemma-cohomology-of-hocolim}.
\end{proof}







\section{Derived limits}
\label{section-derived-limit}

\noindent
In a triangulated category there is a notion of derived limit.

\begin{definition}
\label{definition-derived-limit}
Let $\mathcal{D}$ be a triangulated category.
Let $(K_n, f_n)$ be an inverse system of objects of $\mathcal{D}$.
We say an object $K$ is a {\it derived limit}, or a
{\it homotopy limit} of the system $(K_n)$ if
the product $\prod K_n$ exists and there is a distinguished triangle
$$
K \to \prod K_n \to \prod K_n \to K[1]
$$
where the map $\prod K_n \to \prod K_n$ is given
by $(k_n) \mapsto (k_n - f_{n + 1}(k_{n + 1}))$. If this is the
case, then we sometimes indicate this by the notation $K = R\lim K_n$.
\end{definition}

\noindent
By TR3 a derived limit, if it exists, is unique up to (non-unique)
isomorphism. Moreover, by TR1 a derived limit $R\lim K_n$ exists
as soon as $\prod K_n$ exists. The derived category $D(\textit{Ab})$
of the category of abelian groups is an example of a triangulated category
where all derived limits exist.

\medskip\noindent
The nonuniqueness makes it hard to pin down the derived limit. In
More on Algebra, Lemma \textnormal{[original label: \nolinkurl{more-algebra-lemma-map-into-Rlim}]}
the reader finds an exact sequence
$$
0 \to R^1\lim \Hom(L, K_n[-1]) \to \Hom(L, R\lim K_n)
\to \lim \Hom(L, K_n) \to 0
$$
describing the $\Hom$s into a derived limit in terms of the
usual $\Hom$s.

\begin{lemma}
\label{lemma-products}
Let $\mathcal{A}$ be an abelian category with exact
countable products. Then
\begin{enumerate}
\item $D(\mathcal{A})$ has countable products,
\item countable products $\prod K_i$ in $D(\mathcal{A})$ are obtained by
taking termwise products of any complexes representing the $K_i$, and
\item $H^p(\prod K_i) = \prod H^p(K_i)$.
\end{enumerate}
\end{lemma}

\begin{proof}
Let $K_i^\bullet$ be a complex representing $K_i$ in $D(\mathcal{A})$.
Let $L^\bullet$ be a complex. Suppose given maps
$\alpha_i : L^\bullet \to K_i^\bullet$ in $D(\mathcal{A})$.
This means there exist quasi-isomorphisms $s_i : K_i^\bullet \to M_i^\bullet$
of complexes and maps of complexes $f_i : L^\bullet \to M_i^\bullet$
such that $\alpha_i = s_i^{-1}f_i$. By assumption the map of complexes
$$
s : \prod K_i^\bullet \longrightarrow \prod M_i^\bullet
$$
is a quasi-isomorphism. Hence setting $f = \prod f_i$ we see that
$\alpha = s^{-1}f$ is a map in $D(\mathcal{A})$ whose composition
with the projection $\prod K_i^\bullet \to K_i^\bullet$ is $\alpha_i$.
We omit the verification that $\alpha$ is unique.
\end{proof}

\noindent
The duals of Lemmas \ref{lemma-compute-colimit},
\ref{lemma-colim-hocolim}, and
\ref{lemma-commutes-with-countable-sums}
should be stated here and proved. However, we do not know any applications
of these lemmas for now.

\begin{lemma}
\label{lemma-inverse-limit-bounded-below}
Let $\mathcal{A}$ be an abelian category with countable products and
enough injectives. Let $(K_n)$ be an inverse system of $D^+(\mathcal{A})$.
Then $R\lim K_n$ exists.
\end{lemma}

\begin{proof}
It suffices to show that $\prod K_n$ exists in $D(\mathcal{A})$.
For every $n$ we can represent $K_n$ by a bounded below complex
$I_n^\bullet$ of injectives (Lemma \textnormal{[original label: \nolinkurl{lemma-injective-resolutions-exist}]}).
Then $\prod K_n$ is represented by $\prod I_n^\bullet$, see
Lemma \textnormal{[original label: \nolinkurl{lemma-product-K-injective}]}.
\end{proof}

\begin{remark}
\label{remark-functoriality-derived-limit}
Let $\mathcal{D}$ be a triangulated category. Let $(K_n)$ and $(L_n)$ be
inverse systems of objects of $\mathcal{D}$ with derived limits $K$ and $L$.
Suppose we have a morphism $a : (K_n) \to (L_n)$ of pro-objects, see
Categories, Example \textnormal{[original label: \nolinkurl{categories-example-pro-morphism-inverse-systems}]}.
This means we are given maps $a_n : K_{m(n)} \to L_n$. We may assume
$m(1) < m(2) < m(3) < \ldots$ and that the maps $a_n$ are compatible
with the transition maps in the systems. Then we can consider the maps
$$
a', a'' : \prod K_n \longrightarrow \prod L_n
$$
defined by the rules $a'((k_n)) = (a_n(k_{m(n)}))$ and
$$
a''((k_n)) =
(a_n(k_{m(n)} + f(k_{m(n) + 1}) + \ldots + f(k_{m(n + 1) - 1})))
$$
where each occurence of $f$ denotes a suitable transition map of
the inverse system $(K_n)$. Then the solid diagram
$$
\xymatrix{
K \ar[r] \ar@{..>}[d] & \prod K_n \ar[r] \ar[d]^{a'} &
\prod K_n \ar[d]^{a''} \\
L \ar[r] & \prod L_n \ar[r] & \prod L_n
}
$$
is commutative and by TR3 we obtain the dotted arrow producing
a morphism of distinguished triangles. We warn the reader that
the map $K \to L$ is not unique. We will see later, that if $a$
is a pro-isomorphism, then $K \to L$ is an isomorphism, see
More on Algebra, Lemma \textnormal{[original label: \nolinkurl{more-algebra-lemma-pro-isomorphism-Rlim}]}.
\end{remark}

\begin{remark}
\label{remark-map-into-derived-limit-truncations}
Let $\mathcal{A}$ be an abelian category. Let $K^\bullet$ be a
complex of $\mathcal{A}$. Then $\tau_{\geq -n}K^\bullet$ is an
inverse system of complexes and which in particular determines
an inverse system in $D(\mathcal{A})$. Let us assume that
$R\lim \tau_{\geq -n}K^\bullet$ exists. Then the canonical maps
$c_n : K^\bullet \to \tau_{\geq -n}K^\bullet$ are compatible with the
transition maps of our inverse system. By the defining distinguished
triangle of Definition \ref{definition-derived-limit} and
Lemma \textnormal{[original label: \nolinkurl{lemma-representable-homological}]}
we conclude there exists a morphism
$$
c : K^\bullet \longrightarrow R\lim \tau_{\geq -n}K^\bullet
$$
in $D(\mathcal{A})$ such that the composition of $c$ with the projection
$R\lim \tau_{\geq -n}K^\bullet \to \tau_{\geq -m}K^\bullet$ is equal to $c_m$.
Now the morphism $c$ may not be unique, but we claim that whether or not
$c$ is an isomorphism is independent of the choice of $c$ (and of our choice
of the homotopy limit). Namely, for $i \in \mathbf{Z}$ and for $m > -i$
the composition
$$
H^i(K^\bullet) \xrightarrow{H^i(c)} H^i(R\lim \tau_{\geq -n}K^\bullet)
\to H^i(\tau_{\geq -m}K^\bullet) = H^i(K^\bullet)
$$
is the identity. Hence $H^i(c)$ is an isomorphism if and only if the
second map is an isomorphism. This is independent of $c$ and also
independent of the choice of the homotopy limit (as any two choices
are isomorphic).
\end{remark}

\begin{lemma}
\label{lemma-difficulty-K-injectives}
Let $\mathcal{A}$ be an abelian category with countable products and
enough injectives. Let $K^\bullet$ be a complex. Let $I_n^\bullet$ be
the inverse system of bounded below complexes of injectives produced by
Lemma \textnormal{[original label: \nolinkurl{lemma-special-inverse-system}]}. Then
$I^\bullet = \lim I_n^\bullet$ exists, is K-injective, represents
$R\lim \tau_{\geq -n}K^\bullet$ in $D(\mathcal{A})$,
and the following are equivalent
\begin{enumerate}
\item the map $K^\bullet \to I^\bullet$ (see proof) is a quasi-isomorphism,
\item the map $K^\bullet \to R\lim \tau_{\geq -n}K^\bullet$
of Remark \ref{remark-map-into-derived-limit-truncations}
is an isomorphism in $D(\mathcal{A})$.
\end{enumerate}
\end{lemma}

\begin{proof}
The statement of the lemma makes sense as $R\lim \tau_{\geq -n}K^\bullet$
exists by Lemma \ref{lemma-inverse-limit-bounded-below}.
Each complex $I_n^\bullet$ is K-injective by
Lemma \textnormal{[original label: \nolinkurl{lemma-bounded-below-injectives-K-injective}]}.
Choose direct sum decompositions $I_{n + 1}^p = C_{n + 1}^p \oplus I_n^p$
for all $n \geq 1$. Set $C_1^p = I_1^p$. The complex
$I^\bullet = \lim I_n^\bullet$ exists because we can take
$I^p = \prod_{n \geq 1} C_n^p$. Fix $p \in \mathbf{Z}$.
We claim there is a split short exact sequence
$$
0 \to I^p \to \prod I_n^p \to \prod I_n^p \to 0
$$
of objects of $\mathcal{A}$. Here the first map is given by
the projection maps $I^p \to I_n^p$ and the second map
by $(x_n) \mapsto (x_n - f^p_{n + 1}(x_{n + 1}))$ where
$f^p_n : I_n^p \to I_{n - 1}^p$ are the transition maps.
The splitting comes from the map $\prod I_n^p \to \prod C_n^p = I^p$.
We obtain a termwise split short exact sequence of complexes
$$
0 \to I^\bullet \to \prod I_n^\bullet \to \prod I_n^\bullet \to 0
$$
Hence a corresponding distinguished triangle in $K(\mathcal{A})$
and $D(\mathcal{A})$. By Lemma \textnormal{[original label: \nolinkurl{lemma-product-K-injective}]}
the products are K-injective and represent the corresponding
products in $D(\mathcal{A})$.
It follows that $I^\bullet$ represents $R\lim I_n^\bullet$
(Definition \ref{definition-derived-limit}).
Since $R\lim I_n^\bullet \cong R\lim \tau_{\geq -n}K^\bullet$
as derived limits are defined on the level of the derived
category, we see that $I^\bullet$ represents $R\lim \tau_{\geq -n}K^\bullet$.
Moreover, the complex $I^\bullet$ is K-injective by
Lemma \textnormal{[original label: \nolinkurl{lemma-triangle-K-injective}]}.
By the commutative diagram of Lemma \textnormal{[original label: \nolinkurl{lemma-special-inverse-system}]}
and since $K^i = (\tau_{\geq -n}K^\bullet)^i$
for $n \gg 0$ we see that we get a unique map
$\gamma : K^\bullet \to I^\bullet$ such that the diagrams
$$
\xymatrix{
K^\bullet \ar[r] \ar[d]_\gamma & \tau_{\geq -n} K^\bullet \ar[d] \\
I^\bullet \ar[r] & I_n^\bullet
}
$$
commute. It follows that $\gamma$ is a map of complexes which represents
the map $c : K^\bullet \to R\lim \tau_{\geq -n}K^\bullet$ of
Remark \ref{remark-map-into-derived-limit-truncations}
in $D(\mathcal{A})$. In other words, the diagram
$$
\xymatrix{
K^\bullet \ar[r]_-c \ar[d]_\gamma & R\lim \tau_{\geq -n} K^\bullet
\ar[d]^{\cong} \\
I^\bullet \ar[r]^-{\cong} & R\lim I_n^\bullet
}
$$
is commutative in $D(\mathcal{A})$. The lemma follows.
\end{proof}

\begin{lemma}
\label{lemma-enough-K-injectives-Ab4-star}
Let $\mathcal{A}$ be an abelian category having enough injectives
and exact countable products. Then for every complex
there is a quasi-isomorphism to a K-injective complex.
\end{lemma}

\begin{proof}
By Lemma \ref{lemma-difficulty-K-injectives} it suffices to show that
$K \to R\lim\tau_{\geq -n}K$ is an isomorphism for all $K$ in $D(\mathcal{A})$.
Consider the defining distinguished triangle
$$
R\lim\tau_{\geq -n}K \to
\prod \tau_{\geq -n}K \to
\prod \tau_{\geq -n}K \to
(R\lim\tau_{\geq -n}K)[1]
$$
By Lemma \ref{lemma-products} we have
$$
H^p(\prod \tau_{\geq -n}K) = \prod\nolimits_{p \geq -n} H^p(K)
$$
It follows in a straightforward manner from the long exact cohomology
sequence of the displayed distinguished triangle
that $H^p(R\lim \tau_{\geq -n}K) = H^p(K)$.
\end{proof}




\section{Operations on full subcategories}
\label{section-operate-on-full}

\noindent
Let $\mathcal{T}$ be a triangulated category. We will identify full
subcategories of $\mathcal{T}$ with subsets of $\Ob(\mathcal{T})$.
Given full subcategories $\mathcal{A}, \mathcal{B}, \ldots$ we let
\begin{enumerate}
\item $\mathcal{A}[a, b]$ for $-\infty \leq a \leq b \leq \infty$
be the full subcategory of $\mathcal{T}$ consisting of all objects
$A[-i]$ with $i \in [a, b] \cap \mathbf{Z}$
with $A \in \Ob(\mathcal{A})$ (note the minus sign!),
\item $smd(\mathcal{A})$ be the full subcategory of $\mathcal{T}$
consisting of all objects which are isomorphic to direct summands of objects
of $\mathcal{A}$,
\item $add(\mathcal{A})$ be the full subcategory of $\mathcal{T}$
consisting of all objects which are isomorphic to finite direct sums of objects
of $\mathcal{A}$,
\item $\mathcal{A} \star \mathcal{B}$ be the full subcategory of $\mathcal{T}$
consisting of all objects $X$ of $\mathcal{T}$ which fit into a distinguished
triangle $A \to X \to B$ with $A \in \Ob(\mathcal{A})$ and
$B \in \Ob(\mathcal{B})$,
\item $\mathcal{A}^{\star n} = \mathcal{A} \star \ldots \star \mathcal{A}$
with $n \geq 1$ factors (we will see $\star$ is associative below),
\item $smd(add(\mathcal{A})^{\star n}) =
smd(add(\mathcal{A}) \star \ldots \star add(\mathcal{A}))$
with $n \geq 1$ factors.
\end{enumerate}
If $E$ is an object of $\mathcal{T}$, then we think of $E$ sometimes
also as the full subcategory of $\mathcal{T}$ whose single object is $E$.
Then we can consider things like $add(E[-1, 2])$ and so on and so forth.
We warn the reader that this notation is not universally accepted.

\begin{lemma}
\label{lemma-associativity-star}
Let $\mathcal{T}$ be a triangulated category.
Given full subcategories $\mathcal{A}$, $\mathcal{B}$, $\mathcal{C}$
we have $(\mathcal{A} \star \mathcal{B}) \star \mathcal{C} =
\mathcal{A} \star (\mathcal{B} \star \mathcal{C})$.
\end{lemma}

\begin{proof}
If we have distinguished triangles $A \to X \to B$ and $X \to Y \to C$
then by Axiom TR4 we have distinguished triangles
$A \to Y \to Z$ and $B \to Z \to C$.
\end{proof}

\begin{lemma}
\label{lemma-smd-star}
Let $\mathcal{T}$ be a triangulated category.
Given full subcategories $\mathcal{A}$, $\mathcal{B}$
we have
$smd(\mathcal{A}) \star smd(\mathcal{B}) \subset
smd(\mathcal{A} \star \mathcal{B})$ and
$smd(smd(\mathcal{A}) \star smd(\mathcal{B})) =
smd(\mathcal{A} \star \mathcal{B})$.
\end{lemma}

\begin{proof}
Suppose we have a distinguished triangle $A_1 \to X \to B_1$ where
$A_1 \oplus A_2 \in \Ob(\mathcal{A})$ and $B_1 \oplus B_2 \in \Ob(\mathcal{B})$.
Then we obtain a distinguished triangle
$A_1 \oplus A_2 \to A_2 \oplus X \oplus B_2 \to B_1 \oplus B_2$
which proves that $X$ is in $smd(\mathcal{A} \star \mathcal{B})$.
This proves the inclusion. The equality follows trivially from this.
\end{proof}

\begin{lemma}
\label{lemma-add-star}
Let $\mathcal{T}$ be a triangulated category. Given full subcategories
$\mathcal{A}$, $\mathcal{B}$ the full subcategories
$add(\mathcal{A}) \star add(\mathcal{B})$ and
$smd(add(\mathcal{A}))$ are closed under direct sums.
\end{lemma}

\begin{proof}
Namely, if $A \to X \to B$ and $A' \to X' \to B'$ are distinguished triangles
and $A, A' \in add(\mathcal{A})$ and $B, B' \in add(\mathcal{B})$ then
$A \oplus A' \to X \oplus X' \to B \oplus B'$ is a distinguished triangle
with $A \oplus A' \in add(\mathcal{A})$ and $B \oplus B' \in add(\mathcal{B})$.
The result for $smd(add(\mathcal{A}))$ is trivial.
\end{proof}

\begin{lemma}
\label{lemma-cone-n}
Let $\mathcal{T}$ be a triangulated category. Given a full subcategory
$\mathcal{A}$ for $n \geq 1$ the subcategory
$$
\mathcal{C}_n = smd(add(\mathcal{A})^{\star n}) =
smd(add(\mathcal{A}) \star \ldots \star add(\mathcal{A}))
$$
defined above is a strictly full subcategory of $\mathcal{T}$
closed under direct sums and direct summands and
$\mathcal{C}_{n + m} = smd(\mathcal{C}_n \star \mathcal{C}_m)$
for all $n, m \geq 1$.
\end{lemma}

\begin{proof}
Immediate from Lemmas \ref{lemma-associativity-star}, \ref{lemma-smd-star}, and
\ref{lemma-add-star}.
\end{proof}

\begin{remark}
\label{remark-operations-functor}
Let $F : \mathcal{T} \to \mathcal{T}'$ be an exact functor of triangulated
categories. Given a full subcategory $\mathcal{A}$ of $\mathcal{T}$ we denote
$F(\mathcal{A})$ the full subcategory of $\mathcal{T}'$ whose objects
consists of all objects $F(A)$ with $A \in \Ob(\mathcal{A})$. We have
$$
F(\mathcal{A}[a, b]) = F(\mathcal{A})[a, b]
$$
$$
F(smd(\mathcal{A})) \subset smd(F(\mathcal{A})),
$$
$$
F(add(\mathcal{A})) \subset add(F(\mathcal{A})),
$$
$$
F(\mathcal{A} \star \mathcal{B}) \subset F(\mathcal{A}) \star F(\mathcal{B}),
$$
$$
F(\mathcal{A}^{\star n}) \subset F(\mathcal{A})^{\star n}.
$$
We omit the trivial verifications.
\end{remark}

\begin{remark}
\label{remark-operations-unions}
Let $\mathcal{T}$ be a triangulated category. Given full subcategories
$\mathcal{A}_1 \subset \mathcal{A}_2 \subset \mathcal{A}_3 \subset \ldots$
and $\mathcal{B}$ of $\mathcal{T}$ we have
$$
\left(\bigcup \mathcal{A}_i\right)[a, b] = \bigcup \mathcal{A}_i[a, b]
$$
$$
smd\left(\bigcup \mathcal{A}_i\right) = \bigcup smd(\mathcal{A}_i),
$$
$$
add\left(\bigcup \mathcal{A}_i\right) = \bigcup add(\mathcal{A}_i),
$$
$$
\left(\bigcup \mathcal{A}_i\right) \star \mathcal{B} =
\bigcup \mathcal{A}_i \star \mathcal{B},
$$
$$
\mathcal{B} \star \left(\bigcup \mathcal{A}_i\right) =
\bigcup \mathcal{B} \star \mathcal{A}_i,
$$
$$
\left(\bigcup \mathcal{A}_i\right)^{\star n} =
\bigcup \mathcal{A}_i^{\star n}.
$$
We omit the trivial verifications.
\end{remark}

\begin{lemma}
\label{lemma-in-cone-n}
Let $\mathcal{A}$ be an abelian category. Let $\mathcal{D} = D(\mathcal{A})$.
Let $\mathcal{E} \subset \Ob(\mathcal{A})$ be a subset which we view as
a subset of $\Ob(\mathcal{D})$ also. Let $K$ be an object of $\mathcal{D}$.
\begin{enumerate}
\item Let $b \geq a$ and assume $H^i(K)$ is zero for $i \not \in [a, b]$
and $H^i(K) \in \mathcal{E}$ if $i \in [a, b]$. Then $K$ is in
$smd(add(\mathcal{E}[a, b])^{\star (b - a + 1)})$.
\item Let $b \geq a$ and assume $H^i(K)$ is zero for $i \not \in [a, b]$
and $H^i(K) \in smd(add(\mathcal{E}))$ if $i \in [a, b]$. Then $K$ is in
$smd(add(\mathcal{E}[a, b])^{\star (b - a + 1)})$.
\item Let $b \geq a$ and assume $K$ can be represented by a complex $K^\bullet$
with $K^i = 0$ for $i \not \in [a, b]$ and $K^i \in \mathcal{E}$ for
$i \in [a, b]$. Then $K$ is in
$smd(add(\mathcal{E}[a, b])^{\star (b - a + 1)})$.
\item Let $b \geq a$ and assume $K$ can be represented by a complex $K^\bullet$
with $K^i = 0$ for $i \not \in [a, b]$ and $K^i \in smd(add(\mathcal{E}))$ for
$i \in [a, b]$. Then $K$ is in
$smd(add(\mathcal{E}[a, b])^{\star (b - a + 1)})$.
\end{enumerate}
\end{lemma}

\begin{proof}
We will use Lemma \ref{lemma-cone-n} without further mention.
We will prove (2) which trivially implies (1). We use induction
on $b - a$. If $b - a = 0$, then $K$ is isomorphic to $H^i(K)[-a]$
in $\mathcal{D}$ and the result is immediate. If $b - a > 0$, then
we consider the distinguished triangle
$$
\tau_{\leq b - 1}K^\bullet \to K^\bullet \to K^b[-b]
$$
and we conclude by induction on $b - a$. We omit the proof of (3) and (4).
\end{proof}

\begin{lemma}
\label{lemma-forward-cone-n}
Let $\mathcal{T}$ be a triangulated category. Let
$H : \mathcal{T} \to \mathcal{A}$ be a homological functor
to an abelian category $\mathcal{A}$.
Let $a \leq b$ and $\mathcal{E} \subset \Ob(\mathcal{T})$
be a subset such that $H^i(E) = 0$ for $E \in \mathcal{E}$
and $i \not \in [a, b]$.
Then for $X \in smd(add(\mathcal{E}[-m, m])^{\star n})$
we have $H^i(X) = 0$ for $i \not \in [-m + na, m + nb]$.
\end{lemma}

\begin{proof}
Omitted. Pleasant exercise in the definitions.
\end{proof}





\section{Generators of triangulated categories}
\label{section-generators}

\noindent
In this section we briefly introduce a few of the different notions
of a generator for a triangulated category. Our terminology is
taken from \cite{BvdB} (except that we use ``saturated'' for what
they call ``\'epaisse'', see Definition \textnormal{[original label: \nolinkurl{definition-saturated}]}, and
our definition of $add(\mathcal{A})$ is different).

\medskip\noindent
Let $\mathcal{D}$ be a triangulated category. Let $E$ be an object
of $\mathcal{D}$. Denote $\langle E \rangle_1$ the strictly full subcategory
of $\mathcal{D}$ consisting of objects in $\mathcal{D}$ isomorphic to
direct summands of finite direct sums
$$
\bigoplus\nolimits_{i = 1, \ldots, r} E[n_i]
$$
of shifts of $E$. It is clear that in the notation of
Section \ref{section-operate-on-full} we have
$$
\langle E \rangle_1 = smd(add(E[-\infty, \infty]))
$$
For $n > 1$ let $\langle E \rangle_n$ denote the full
subcategory of $\mathcal{D}$ consisting of objects of $\mathcal{D}$
isomorphic to direct summands of objects $X$ which fit into a distinguished
triangle
$$
A \to X \to B \to A[1]
$$
where $A$ is an object of $\langle E \rangle_1$ and $B$ an object of
$\langle E \rangle_{n - 1}$. In the notation of
Section \ref{section-operate-on-full} we have
$$
\langle E \rangle_n = smd(\langle E \rangle_1 \star \langle E \rangle_{n - 1})
$$
Each of the categories $\langle E \rangle_n$ is a strictly full additive
(by Lemma \ref{lemma-add-star}) subcategory of $\mathcal{D}$ preserved
under shifts and under taking summands. But, $\langle E \rangle_n$ is not
necessarily closed under ``taking cones'' or ``extensions'', hence not
necessarily a triangulated subcategory. This will be true for the
subcategory
$$
\langle E \rangle = \bigcup\nolimits_n \langle E \rangle_n
$$
as will be shown in the lemmas below.

\begin{lemma}
\label{lemma-generated-by-E-explicit}
Let $\mathcal{T}$ be a triangulated category. Let $E$ be an object
of $\mathcal{T}$. For $n \geq 1$ we have
$$
\langle E \rangle_n =
smd(\langle E \rangle_1 \star \ldots \star \langle E \rangle_1) =
smd({\langle E \rangle_1}^{\star n}) =
\bigcup\nolimits_{m \geq 1} smd(add(E[-m, m])^{\star n})
$$
For $n, n' \geq 1$ we have $\langle E \rangle_{n + n'} =
smd(\langle E \rangle_n \star \langle E \rangle_{n'})$.
\end{lemma}

\begin{proof}
The left equality in the displayed formula follows from
Lemmas \ref{lemma-associativity-star} and \ref{lemma-smd-star}
and induction. The middle equality is a matter of notation.
Since $\langle E \rangle_1 = smd(add(E[-\infty, \infty])])$
and since $E[-\infty, \infty] = \bigcup_{m \geq 1} E[-m, m]$
we see from Remark \ref{remark-operations-unions} and
Lemma \ref{lemma-smd-star} that we get the equality on the right.
Then the final statement follows from the remark and the
corresponding statement of Lemma \ref{lemma-cone-n}.
\end{proof}

\begin{lemma}
\label{lemma-find-smallest-containing-E}
Let $\mathcal{D}$ be a triangulated category. Let $E$ be an object
of $\mathcal{D}$. The subcategory
$$
\langle E \rangle = \bigcup\nolimits_n \langle E \rangle_n
= \bigcup\nolimits_{n, m \geq 1} smd(add(E[-m, m])^{\star n})
$$
is a strictly full, saturated, triangulated subcategory of $\mathcal{D}$
and it is the smallest such subcategory of $\mathcal{D}$ containing
the object $E$.
\end{lemma}

\begin{proof}
The equality on the right follows from
Lemma \ref{lemma-generated-by-E-explicit}.
It is clear that $\langle E \rangle = \bigcup \langle E \rangle_n$
contains $E$, is preserved under shifts, direct sums, direct summands.
If $A \in \langle E \rangle_a$ and $B \in \langle E \rangle_b$
and if $A \to X \to B \to A[1]$ is a distinguished triangle, then
$X \in \langle E \rangle_{a + b}$ by Lemma \ref{lemma-generated-by-E-explicit}.
Hence $\bigcup \langle E \rangle_n$ is also preserved under extensions
and it follows that it is a triangulated subcategory. 

\medskip\noindent
Finally, let $\mathcal{D}' \subset \mathcal{D}$ be a 
strictly full, saturated, triangulated subcategory of $\mathcal{D}$
containing $E$. Then
$\mathcal{D}'[-\infty, \infty] \subset \mathcal{D}'$,
$add(\mathcal{D}) \subset \mathcal{D}'$,
$smd(\mathcal{D}') \subset \mathcal{D}'$, and
$\mathcal{D}' \star \mathcal{D}' \subset \mathcal{D}'$.
In other words, all the operations we used to construct
$\langle E \rangle$ out of $E$ preserve $\mathcal{D}'$.
Hence $\langle E \rangle \subset \mathcal{D}'$ and this
finishes the proof.
\end{proof}

\begin{definition}
\label{definition-generators}
Let $\mathcal{D}$ be a triangulated category. Let $E$ be an object
of $\mathcal{D}$.
\begin{enumerate}
\item We say $E$ is a {\it classical generator} of $\mathcal{D}$
if the smallest strictly full, saturated, triangulated subcategory
of $\mathcal{D}$ containing $E$ is equal to $\mathcal{D}$, in
other words, if $\langle E \rangle = \mathcal{D}$.
\item We say $E$ is a {\it strong generator} of $\mathcal{D}$
if $\langle E \rangle_n = \mathcal{D}$ for some $n \geq 1$.
\item We say $E$ is a {\it weak generator} or a {\it generator}
of $\mathcal{D}$
if for any nonzero object $K$ of $\mathcal{D}$ there exists
an integer $n$ and a nonzero map $E \to K[n]$.
\end{enumerate}
\end{definition}

\noindent
This definition can be generalized to the case of a family of objects.

\begin{lemma}
\label{lemma-right-orthogonal}
Let $\mathcal{D}$ be a triangulated category. Let $E, K$ be objects
of $\mathcal{D}$. The following are equivalent
\begin{enumerate}
\item $\Hom(E, K[i]) = 0$ for all $i \in \mathbf{Z}$,
\item $\Hom(E', K) = 0$ for all $E' \in \langle E \rangle$.
\end{enumerate}
\end{lemma}

\begin{proof}
The implication (2) $\Rightarrow$ (1) is immediate. Conversely,
assume (1). Then $\Hom(X, K) = 0$ for all $X$ in $\langle E \rangle_1$.
Arguing by induction on $n$ and using
Lemma \textnormal{[original label: \nolinkurl{lemma-representable-homological}]}
we see that $\Hom(X, K) = 0$ for all $X$
in $\langle E \rangle_n$.
\end{proof}

\begin{lemma}
\label{lemma-classical-generator-generator}
Let $\mathcal{D}$ be a triangulated category. Let $E$ be an object
of $\mathcal{D}$. If $E$ is a classical generator of $\mathcal{D}$,
then $E$ is a generator.
\end{lemma}

\begin{proof}
Assume $E$ is a classical generator. Let $K$ be an object of $\mathcal{D}$
such that $\Hom(E, K[i]) = 0$ for all $i \in \mathbf{Z}$. By
Lemma \ref{lemma-right-orthogonal}
$\Hom(E', K) = 0$ for all $E'$ in $\langle E \rangle$. However, since
$\mathcal{D} = \langle E \rangle$ we conclude that $\text{id}_K = 0$,
i.e., $K = 0$.
\end{proof}

\begin{lemma}
\label{lemma-classical-generator-strong-generator}
Let $\mathcal{D}$ be a triangulated category which has a strong generator.
Let $E$ be an object of $\mathcal{D}$. If $E$ is a classical generator of
$\mathcal{D}$, then $E$ is a strong generator.
\end{lemma}

\begin{proof}
Let $E'$ be an object of $\mathcal{D}$ such that
$\mathcal{D} = \langle E' \rangle_n$. Since
$\mathcal{D} = \langle E \rangle$ we see that $E' \in \langle E \rangle_m$
for some $m \geq 1$ by Lemma \ref{lemma-find-smallest-containing-E}.
Then $\langle E' \rangle_1 \subset \langle E \rangle_m$ hence
$$
\mathcal{D} =
\langle E' \rangle_n = smd(
\langle E' \rangle_1 \star \ldots \star \langle E' \rangle_1)
\subset
smd(
\langle E \rangle_m \star \ldots \star \langle E \rangle_m)
=
\langle E \rangle_{nm}
$$
as desired. Here we used Lemma \ref{lemma-generated-by-E-explicit}.
\end{proof}

\begin{remark}
\label{remark-check-on-generator}
Let $\mathcal{D}$ be a triangulated category. Let $E$ be an object
of $\mathcal{D}$. Let $T$ be a property of objects of $\mathcal{D}$.
Suppose that
\begin{enumerate}
\item if $K_i \in \mathcal{D}$, $i = 1, \ldots, r$ with
$T(K_i)$ for $i = 1, \ldots, r$, then $T(\bigoplus K_i)$,
\item if $K \to L \to M \to K[1]$ is a distinguished triangle and
$T$ holds for two, then $T$ holds for the third object,
\item if $T(K \oplus L)$ then $T(K)$ and $T(L)$, and
\item $T(E[n])$ holds for all $n$.
\end{enumerate}
Then $T$ holds for all objects of $\langle E \rangle$.
\end{remark}












\section{Compact objects}
\label{section-compact}

\noindent
Here is the definition.

\begin{definition}
\label{definition-compact-object}
Let $\mathcal{D}$ be an additive category with arbitrary direct
sums. A {\it compact object} of $\mathcal{D}$ is an object $K$
such that the map
$$
\bigoplus\nolimits_{i \in I} \Hom_{\mathcal{D}}(K, E_i)
\longrightarrow
\Hom_{\mathcal{D}}(K, \bigoplus\nolimits_{i \in I} E_i)
$$
is bijective for any set $I$ and objects
$E_i \in \Ob(\mathcal{D})$ parametrized by $i \in I$.
\end{definition}

\noindent
This notion turns out to be very useful in algebraic geometry.
It is an intrinsic condition on objects that forces the objects
to be, well, compact.

\begin{lemma}
\label{lemma-compact-objects-subcategory}
Let $\mathcal{D}$ be a (pre-)triangulated category with direct sums.
Then the compact objects of $\mathcal{D}$ form the objects of a
Karoubian, saturated, strictly full, (pre-)triangulated subcategory
$\mathcal{D}_c$ of $\mathcal{D}$.
\end{lemma}

\begin{proof}
Let $(X, Y, Z, f, g, h)$ be a distinguished triangle of $\mathcal{D}$
with $X$ and $Y$ compact. Then it follows from
Lemma \textnormal{[original label: \nolinkurl{lemma-representable-homological}]}
and the five lemma
(Homology, Lemma \textnormal{[original label: \nolinkurl{homology-lemma-five-lemma}]})
that $Z$ is a compact object too. It is clear that if $X \oplus Y$
is compact, then $X$, $Y$ are compact objects too. Hence
$\mathcal{D}_c$ is a saturated triangulated subcategory.
Since $\mathcal{D}$ is Karoubian by
Lemma \textnormal{[original label: \nolinkurl{lemma-projectors-have-images-triangulated}]}
we conclude that the same is true for $\mathcal{D}_c$.
\end{proof}

\begin{lemma}
\label{lemma-write-as-colimit}
Let $\mathcal{D}$ be a triangulated category with direct sums.
Let $E_i$, $i \in I$ be a family of compact objects of $\mathcal{D}$
such that $\bigoplus E_i$ generates $\mathcal{D}$.
Then every object $X$ of $\mathcal{D}$ can be written as
$$
X = \text{hocolim} X_n
$$
where $X_1$ is a direct sum of shifts of the $E_i$ and each transition
morphism fits into a distinguished triangle
$Y_n \to X_n \to X_{n + 1} \to Y_n[1]$
where $Y_n$ is a direct sum of shifts of the $E_i$.
\end{lemma}

\begin{proof}
Set $X_1 = \bigoplus_{(i, m, \varphi)} E_i[m]$ where the direct sum is over
all triples $(i, m, \varphi)$ such that $i \in I$, $m \in \mathbf{Z}$
and $\varphi : E_i[m] \to X$. Then $X_1$ comes equipped with a canonical
morphism $X_1 \to X$. Given $X_n \to X$ we set
$Y_n = \bigoplus_{(i, m, \varphi)} E_i[m]$ where the direct sum is over
all triples $(i, m, \varphi)$ such that $i \in I$, $m \in \mathbf{Z}$, and
$\varphi : E_i[m] \to X_n$ is a morphism such that $E_i[m] \to X_n \to X$
is zero. Choose a distinguished triangle
$Y_n \to X_n \to X_{n + 1} \to Y_n[1]$
and let $X_{n + 1} \to X$ be any morphism such that $X_n \to X_{n + 1} \to X$
is the given one; such a morphism exists by our choice of $Y_n$.
We obtain a morphism $\text{hocolim} X_n \to X$ by the construction
of our maps $X_n \to X$. Choose a distinguished triangle
$$
C \to \text{hocolim} X_n \to X \to C[1]
$$
Let $E_i[m] \to C$ be a morphism. Since $E_i$ is compact,
Lemma \ref{lemma-commutes-with-countable-sums}
implies the composition $E_i[m] \to C \to \text{hocolim} X_n$ factors through
$X_n$ for some $n$, say by $E_i[m] \to X_n$. Then the
construction of $Y_n$ shows that the composition
$E_i[m] \to X_n \to X_{n + 1}$ is zero. In other words, the composition
$E_i[m] \to C \to \text{hocolim} X_n$ is zero. This means that our
morphism $E_i[m] \to C$ comes from a morphism $E_i[m] \to X[-1]$.
The construction of $X_1$ then shows that such morphism lifts to
$\text{hocolim} X_n$ and we conclude that our morphism $E_i[m] \to C$
is zero. The assumption that $\bigoplus E_i$ generates $\mathcal{D}$
implies that $C$ is zero and the proof is done.
\end{proof}

\begin{lemma}
\label{lemma-factor-through}
With assumptions and notation as in Lemma \ref{lemma-write-as-colimit}.
If $C$ is a compact object and $C \to X_n$ is a morphism, then
there is a factorization $C \to E \to X_n$ where
$E$ is an object of $\langle E_{i_1} \oplus \ldots \oplus E_{i_t} \rangle$
for some $i_1, \ldots, i_t \in I$.
\end{lemma}

\begin{proof}
We prove this by induction on $n$. The base case $n = 1$ is clear.
If $n > 1$ consider the composition $C \to X_n \to Y_{n - 1}[1]$.
This can be factored through some $E'[1] \to Y_{n - 1}[1]$ where
$E'$ is a finite direct sum of shifts of the $E_i$. Let $I' \subset I$
be the finite set of indices that occur in this direct sum. Thus we obtain
$$
\xymatrix{
E' \ar[r] \ar[d] &
C' \ar[r] \ar[d] &
C \ar[r] \ar[d] &
E'[1] \ar[d] \\
Y_{n - 1} \ar[r] &
X_{n - 1} \ar[r] &
X_n \ar[r] &
Y_{n - 1}[1]
}
$$
By induction the morphism $C' \to X_{n - 1}$ factors through
$E'' \to X_{n - 1}$ with $E''$ an object of
$\langle \bigoplus_{i \in I''} E_i \rangle$
for some finite subset $I'' \subset I$. Choose a distinguished
triangle
$$
E' \to E'' \to E \to E'[1]
$$
then $E$ is an object of $\langle \bigoplus_{i \in I' \cup I''} E_i \rangle$.
By construction and the axioms of a triangulated category we can choose
morphisms $C \to E$ and a morphism $E \to X_n$ fitting into morphisms
of triangles $(E', C', C) \to (E', E'', E)$ and
$(E', E'', E) \to (Y_{n - 1}, X_{n - 1}, X_n)$. The composition
$C \to E \to X_n$ may not equal the given morphism $C \to X_n$, but
the compositions into $Y_{n - 1}$ are equal. Let $C \to X_{n - 1}$
be a morphism that lifts the difference. By induction assumption we
can factor this through a morphism $E''' \to X_{n - 1}$ with
$E''$ an object of $\langle \bigoplus_{i \in I'''} E_i \rangle$
for some finite subset $I' \subset I$. Thus we see that we get
a solution on considering $E \oplus E''' \to X_n$ because
$E \oplus E'''$ is an object of
$\langle \bigoplus_{i \in I' \cup I'' \cup I'''} E_i \rangle$.
\end{proof}

\begin{definition}
\label{definition-compactly-generated}
Let $\mathcal{D}$ be a triangulated category with arbitrary direct
sums. We say $\mathcal{D}$ is {\it compactly generated} if
there exists a set $E_i$, $i \in I$ of compact objects such that
$\bigoplus E_i$ generates $\mathcal{D}$.
\end{definition}

\noindent
The following proposition clarifies the relationship between
classical generators and weak generators.

\begin{proposition}
\label{proposition-generator-versus-classical-generator}
Let $\mathcal{D}$ be a triangulated category with direct sums.
Let $E$ be a compact object of $\mathcal{D}$.
The following are equivalent
\begin{enumerate}
\item $E$ is a classical generator for $\mathcal{D}_c$ and
$\mathcal{D}$ is compactly generated, and
\item $E$ is a generator for $\mathcal{D}$.
\end{enumerate}
\end{proposition}

\begin{proof}
If $E$ is a classical generator for $\mathcal{D}_c$, then
$\mathcal{D}_c = \langle E \rangle$. It follows formally
from the assumption that $\mathcal{D}$ is compactly generated
and Lemma \ref{lemma-right-orthogonal} that $E$ is a generator
for $\mathcal{D}$.

\medskip\noindent
The converse is more interesting. Assume that $E$ is a generator
for $\mathcal{D}$. Let $X$ be a compact object of $\mathcal{D}$.
Apply Lemma \ref{lemma-write-as-colimit} with $I = \{1\}$ and
$E_1 = E$ to write
$$
X = \text{hocolim} X_n
$$
as in the lemma. Since $X$ is compact we
find that $X \to \text{hocolim} X_n$ factors through $X_n$ for
some $n$ (Lemma \ref{lemma-commutes-with-countable-sums}).
Thus $X$ is a direct summand of $X_n$.
By Lemma \ref{lemma-factor-through} we see that $X$ is an
object of $\langle E \rangle$ and the lemma is proven.
\end{proof}











\section{Brown representability}
\label{section-brown}

\noindent
A reference for the material in this section is \cite{Neeman-Grothendieck}.

\begin{lemma}
\label{lemma-brown}
\begin{reference}
\cite[Theorem 3.1]{Neeman-Grothendieck}.
\end{reference}
Let $\mathcal{D}$ be a triangulated category with direct sums which is
compactly generated. Let $H : \mathcal{D} \to \textit{Ab}$ be a contravariant
cohomological functor which transforms direct sums into products.
Then $H$ is representable.
\end{lemma}

\begin{proof}
Let $E_i$, $i \in I$ be a set of compact objects such that
$\bigoplus_{i \in I} E_i$ generates $\mathcal{D}$. We may and do assume
that the set of objects $\{E_i\}$ is preserved under shifts. Consider pairs
$(i, a)$ where $i \in I$ and $a \in H(E_i)$ and set
$$
X_1 = \bigoplus\nolimits_{(i, a)} E_i
$$
Since $H(X_1) = \prod_{(i, a)} H(E_i)$ we see that $(a)_{(i, a)}$
defines an element $a_1 \in H(X_1)$. Set $H_1 = \Hom_\mathcal{D}(- , X_1)$.
By Yoneda's lemma (Categories, Lemma \textnormal{[original label: \nolinkurl{categories-lemma-yoneda}]})
the element $a_1$ defines a natural transformation $H_1 \to H$.

\medskip\noindent
We are going to inductively construct $X_n$ and transformations
$a_n : H_n \to H$ where $H_n = \Hom_\mathcal{D}(-, X_n)$.
Namely, we apply the procedure
above to the functor $\Ker(H_n \to H)$ to get an object
$$
K_{n + 1} = \bigoplus\nolimits_{(i, k),\ k \in \Ker(H_n(E_i) \to H(E_i))} E_i
$$
and a transformation $\Hom_\mathcal{D}(-, K_{n + 1}) \to \Ker(H_n \to H)$.
By Yoneda's lemma the composition $\Hom_\mathcal{D}(-, K_{n + 1}) \to H_n$
gives a morphism $K_{n + 1} \to X_n$. We choose
a distinguished triangle
$$
K_{n + 1} \to X_n \to X_{n + 1} \to K_{n + 1}[1]
$$
in $\mathcal{D}$. The element $a_n \in H(X_n)$ maps to zero
in $H(K_{n + 1})$ by construction. Since $H$ is cohomological
we can lift it to an element $a_{n + 1} \in H(X_{n + 1})$.

\medskip\noindent
We claim that $X = \text{hocolim} X_n$ represents $H$. Applying $H$
to the defining distinguished triangle
$$
\bigoplus X_n \to
\bigoplus X_n \to X \to
\bigoplus X_n[1]
$$
we obtain an exact sequence
$$
\prod H(X_n) \leftarrow 
\prod H(X_n) \leftarrow 
H(X)
$$
Thus there exists an element $a \in H(X)$ mapping to $(a_n)$
in $\prod H(X_n)$. Hence a natural transformation
$\Hom_\mathcal{D}(- , X) \to H$ such that
$$
\Hom_\mathcal{D}(-, X_1) \to
\Hom_\mathcal{D}(-, X_2) \to
\Hom_\mathcal{D}(-, X_3) \to \ldots \to
\Hom_\mathcal{D}(-, X) \to H
$$
commutes. For each $i$ the map $\Hom_\mathcal{D}(E_i, X) \to H(E_i)$
is surjective, by construction of $X_1$. On the other hand, by construction
of $X_n \to X_{n + 1}$ the kernel of $\Hom_\mathcal{D}(E_i, X_n) \to H(E_i)$
is killed by the map
$\Hom_\mathcal{D}(E_i, X_n) \to \Hom_\mathcal{D}(E_i, X_{n + 1})$.
Since
$$
\Hom_\mathcal{D}(E_i, X) = \colim \Hom_\mathcal{D}(E_i, X_n)
$$
by Lemma \ref{lemma-commutes-with-countable-sums}
we see that $\Hom_\mathcal{D}(E_i, X) \to H(E_i)$ is injective.

\medskip\noindent
To finish the proof, consider the subcategory
$$
\mathcal{D}' =
\{Y \in \Ob(\mathcal{D}) \mid \Hom_\mathcal{D}(Y[n], X) \to H(Y[n])
\text{ is an isomorphism for all }n\}
$$
As $\Hom_\mathcal{D}(-, X) \to H$ is a transformation between
cohomological functors,
the subcategory $\mathcal{D}'$ is a strictly full, saturated, triangulated
subcategory of $\mathcal{D}$ (details omitted; see proof of
Lemma \textnormal{[original label: \nolinkurl{lemma-homological-functor-kernel}]}). Moreover, as both
$H$ and $\Hom_\mathcal{D}(-, X)$ transform direct sums into products,
we see that direct sums of objects of $\mathcal{D}'$ are in $\mathcal{D}'$.
Thus derived colimits of objects of $\mathcal{D}'$ are in $\mathcal{D}'$.
Since $\{E_i\}$ is preserved under shifts, we see that $E_i$
is an object of $\mathcal{D}'$ for all $i$. It follows from
Lemma \ref{lemma-write-as-colimit} that $\mathcal{D}' = \mathcal{D}$
and the proof is complete.
\end{proof}

\begin{proposition}
\label{proposition-brown}
\begin{reference}
\cite[Theorem 4.1]{Neeman-Grothendieck}.
\end{reference}
Let $\mathcal{D}$ be a triangulated category with direct sums which is
compactly generated. Let $F : \mathcal{D} \to \mathcal{D}'$ be an
exact functor of triangulated categories which transforms direct sums
into direct sums. Then $F$ has an exact right adjoint.
\end{proposition}

\begin{proof}
For an object $Y$ of $\mathcal{D}'$ consider the contravariant functor
$$
\mathcal{D} \to \textit{Ab},\quad W \mapsto \Hom_{\mathcal{D}'}(F(W), Y)
$$
This is a cohomological functor as $F$ is exact and transforms direct sums
into products as $F$ transforms direct sums into direct sums. Thus by
Lemma \ref{lemma-brown} we find an object $X$ of $\mathcal{D}$ such that
$\Hom_\mathcal{D}(W, X) = \Hom_{\mathcal{D}'}(F(W), Y)$.
The existence of the adjoint follows from
Categories, Lemma \textnormal{[original label: \nolinkurl{categories-lemma-adjoint-exists}]}.
Exactness follows from Lemma \textnormal{[original label: \nolinkurl{lemma-adjoint-is-exact}]}.
\end{proof}









\section{Brown representability, bis}
\label{section-brown-bis}

\noindent
In this section we explain a version of Brown representability
for triangulated categories which have a suitable set of generators;
for other versions, please see \cite{Franke}, \cite{Neeman}, and \cite{Krause}.


\begin{lemma}
\label{lemma-brown-bis}
\begin{reference}
Weak version of \cite[Theorem A]{Krause}
\end{reference}
Let $\mathcal{D}$ be a triangulated category with direct sums.
Suppose given a set $\mathcal{E}$ of objects of $\mathcal{D}$ such that
\begin{enumerate}
\item if $X$ is a nonzero object of $\mathcal{D}$, then there exists
an $E \in \mathcal{E}$ and a nonzero map $E \to X$, and
\item given objects $X_n$, $n \in \mathbf{N}$ of $\mathcal{D}$,
$E \in \mathcal{E}$, and $\alpha : E \to \bigoplus X_n$,
there exist $E_n \in \mathcal{E}$ and $\beta_n : E_n \to X_n$ and a morphism
$\gamma : E \to \bigoplus E_n$ such that
$\alpha = (\bigoplus \beta_n) \circ \gamma$.
\end{enumerate}
Let $H : \mathcal{D} \to \textit{Ab}$ be a contravariant cohomological functor
which transforms direct sums into products. Then $H$ is representable.
\end{lemma}



\begin{proof}
This proof is very similar to the proof of Lemma \ref{lemma-brown}.
We may replace $\mathcal{E}$ by
$\bigcup_{i \in \mathbf{Z}} \mathcal{E}[i]$ and assume
that $\mathcal{E}$ is preserved by shifts.
Consider pairs $(E, a)$ where $E \in \mathcal{E}$ and $a \in H(E)$ and set
$$
X_1 = \bigoplus\nolimits_{(E, a)} E
$$
Since $H(X_1) = \prod_{(E, a)} H(E)$ we see that $(a)_{(E, a)}$
defines an element $a_1 \in H(X_1)$. Set $H_1 = \Hom_\mathcal{D}(- , X_1)$.
By Yoneda's lemma (Categories, Lemma \textnormal{[original label: \nolinkurl{categories-lemma-yoneda}]})
the element $a_1$ defines a natural transformation $H_1 \to H$.

\medskip\noindent
We are going to inductively construct $X_n$ and transformations
$a_n : H_n \to H$ where $H_n = \Hom_\mathcal{D}(-, X_n)$.
Namely, we apply the procedure
above to the functor $\Ker(H_n \to H)$ to get an object
$$
K_{n + 1} = \bigoplus\nolimits_{(E, k),\ k \in \Ker(H_n(E) \to H(E))} E
$$
and a transformation $\Hom_\mathcal{D}(-, K_{n + 1}) \to \Ker(H_n \to H)$.
By Yoneda's lemma the composition $\Hom_\mathcal{D}(-, K_{n + 1}) \to H_n$
gives a morphism $K_{n + 1} \to X_n$. We choose
a distinguished triangle
$$
K_{n + 1} \to X_n \to X_{n + 1} \to K_{n + 1}[1]
$$
in $\mathcal{D}$. The element $a_n \in H(X_n)$ maps to zero
in $H(K_{n + 1})$ by construction. Since $H$ is cohomological
we can lift it to an element $a_{n + 1} \in H(X_{n + 1})$.

\medskip\noindent
Set $X = \text{hocolim} X_n$. Applying $H$
to the defining distinguished triangle
$$
\bigoplus X_n \to
\bigoplus X_n \to X \to
\bigoplus X_n[1]
$$
we obtain an exact sequence
$$
\prod H(X_n) \leftarrow
\prod H(X_n) \leftarrow
H(X)
$$
Thus there exists an element $a \in H(X)$ mapping to $(a_n)$
in $\prod H(X_n)$. Hence a natural transformation
$\Hom_\mathcal{D}(- , X) \to H$ such that
$$
\Hom_\mathcal{D}(-, X_1) \to
\Hom_\mathcal{D}(-, X_2) \to
\Hom_\mathcal{D}(-, X_3) \to \ldots \to
\Hom_\mathcal{D}(-, X) \to H
$$
commutes. We claim that $\Hom_\mathcal{D}(-, X) \to H(-)$ is an isomorphism.

\medskip\noindent
Let $E \in \mathcal{E}$. Let us show that
$$
\Hom_\mathcal{D}(E, \bigoplus X_n) \to  \Hom_\mathcal{D}(E, \bigoplus X_n)
$$
is injective. Namely, let $\alpha : E \to \bigoplus X_n$. Then
by assumption (2) we obtain a factorization
$\alpha = (\bigoplus \beta_n) \circ \gamma$.
Since $E_n \to X_n \to X_{n + 1}$ is zero by construction, we see that
the composition $\bigoplus E_n \to \bigoplus X_n \to \bigoplus X_n$
is equal to $\bigoplus \beta_n$. Hence also the composition
$E \to \bigoplus X_n \to \bigoplus X_n$ is equal to $\alpha$.
This proves the stated injectivity and hence also
$$
\Hom_\mathcal{D}(E, \bigoplus X_n[1]) \to \Hom_\mathcal{D}(E, \bigoplus X_n[1])
$$
is injective. It follows that we have an exact sequence
$$
\Hom_\mathcal{D}(E, \bigoplus X_n) \to
\Hom_\mathcal{D}(E, \bigoplus X_n) \to
\Hom_\mathcal{D}(E, X) \to 0
$$
for all $E \in \mathcal{E}$.

\medskip\noindent
Let $E \in \mathcal{E}$ and let $f : E \to X$ be a morphism.
By the previous paragraph, we may choose $\alpha : E \to \bigoplus X_n$
lifting $f$. Then by assumption (2) we obtain a factorization
$\alpha = (\bigoplus \beta_n) \circ \gamma$.
For each $n$ there is a morphism $\delta_n : E_n \to X_1$ such that
$\delta_n$ and $\beta_n$ map to the same element of $H(E_n)$. Then
the compositions
$$
E_n \to X_n \to X_{n + 1}
\quad\text{and}\quad
E_n \to X_1 \to X_{n + 1}
$$
are equal by construction of $X_n \to X_{n + 1}$. It follows that
$$
\bigoplus E_n \to \bigoplus X_n \to X
\quad\text{and}\quad
\bigoplus E_n \to \bigoplus X_1 \to X
$$
are the same too. Observing that $\bigoplus X_1 \to X$ factors
as $\bigoplus X_1 \to X_1 \to X$, we conclude that
$$
\Hom_\mathcal{D}(E, X_1) \to \Hom_\mathcal{D}(E, X)
$$
is surjective. Since by construction the map
$\Hom_\mathcal{D}(E, X_1) \to H(E)$ is surjective and by
construction the kernel of this map is annihilated by
$\Hom_\mathcal{D}(E, X_1) \to \Hom_\mathcal{D}(E, X)$
we conclude that $\Hom_\mathcal{D}(E, X) \to H(E)$ is
a bijection for all $E \in \mathcal{E}$.

\medskip\noindent
To finish the proof, consider the subcategory
$$
\mathcal{D}' =
\{Y \in \Ob(\mathcal{D}) \mid \Hom_\mathcal{D}(Y[n], X) \to H(Y[n])
\text{ is an isomorphism for all }n\}
$$
As $\Hom_\mathcal{D}(-, X) \to H$ is a transformation between
cohomological functors,
the subcategory $\mathcal{D}'$ is a strictly full, saturated, triangulated
subcategory of $\mathcal{D}$ (details omitted; see proof of
Lemma \textnormal{[original label: \nolinkurl{lemma-homological-functor-kernel}]}). Moreover, as both
$H$ and $\Hom_\mathcal{D}(-, X)$ transform direct sums into products,
we see that direct sums of objects of $\mathcal{D}'$ are in $\mathcal{D}'$.
Thus derived colimits of objects of $\mathcal{D}'$ are in $\mathcal{D}'$.
Since $\mathcal{E}$ is preserved by shifts, we conclude that
$\mathcal{E} \subset \Ob(\mathcal{D}')$ by the result of the
previous paragraph. To finish the proof we have to show that
$\mathcal{D}' = \mathcal{D}$.

\medskip\noindent
Let $Y$ be an object of $\mathcal{D}$ and set $H(-) = \Hom_\mathcal{D}(-, Y)$.
Then $H$ is a cohomological functor which transforms direct sums into products.
By the construction in the first part of the proof we obtain a morphism
$\colim X_n = X \to Y$ such that
$\Hom_\mathcal{D}(E, X) \to \Hom_\mathcal{D}(E, Y)$ is bijective for all
$E \in \mathcal{E}$. Then assumption (1) tells us that $X \to Y$ is an
isomorphism! On the other hand, by construction $X_1, X_2, \ldots$
are in $\mathcal{D}'$ and so is $X$. Thus $Y \in \mathcal{D}'$ and
the proof is complete.
\end{proof}


\begin{proposition}
\label{proposition-brown-bis}
Let $\mathcal{D}$ be a triangulated category with direct sums.
Assume there exists a set $\mathcal{E}$ of objects of $\mathcal{D}$
satisfying conditions (1) and (2) of Lemma \ref{lemma-brown-bis}.
Let $F : \mathcal{D} \to \mathcal{D}'$ be an exact functor of triangulated
categories which transforms direct sums into direct sums.
Then $F$ has an exact right adjoint.
\end{proposition}

\begin{proof}
For an object $Y$ of $\mathcal{D}'$ consider the contravariant functor
$$
\mathcal{D} \to \textit{Ab},\quad W \mapsto \Hom_{\mathcal{D}'}(F(W), Y)
$$
This is a cohomological functor as $F$ is exact and transforms direct sums
into products as $F$ transforms direct sums into direct sums. Thus by
Lemma \ref{lemma-brown-bis} we find an object $X$ of $\mathcal{D}$ such that
$\Hom_\mathcal{D}(W, X) = \Hom_{\mathcal{D}'}(F(W), Y)$.
The existence of the adjoint follows from
Categories, Lemma \textnormal{[original label: \nolinkurl{categories-lemma-adjoint-exists}]}.
Exactness follows from Lemma \textnormal{[original label: \nolinkurl{lemma-adjoint-is-exact}]}.
\end{proof}










\bibliographystyle{plain}
\bibliography{my}
\end{document}
